\chapter{Discussion}
\label{ch:discussion}

Every limitation below is a deliberate scope choice or a measured
finding, not an oversight discovered in retrospect; several are results
in their own right. This chapter names each threat to validity in the
standard vocabulary and states where the project mitigated it, and where
it simply lives with it.

\paragraph{Look-ahead bias.} The report contains a full worked example
of this threat rather than merely a disclaimer: the energy track's first
real-data result (\cref{sec:en-artifact}) was look-ahead bias in two of
its three layers, and the diagnostic arc that dismantled it is the
chapter. On the equity track one mild instance is disclosed and
retained: the static correlation graph is built from data up to the
split date, so early \emph{training} snapshots see a graph estimated
partly from their own future. The out-of-sample window is unaffected ---
the graph strictly predates it --- and the clean alternative, the
dynamic graph, was rejected because it performs worse
(\cref{sec:eq-ablation}). Disclosure was chosen over a cosmetic fix that
costs real out-of-sample performance.

\paragraph{Survivorship bias.} The equity universe is roughly fifty
hand-picked, currently listed US large caps fetched from a free vendor.
Names that failed or delisted over 2021--2026 are absent, which inflates
the absolute level of every Sharpe in \cref{ch:equity}. The core claim
is deliberately insulated: the attention A/B compares models on the
\emph{same} biased universe, so the bias raises both sides of the
comparison rather than manufacturing the gap. Absolute levels should be
read as optimistic; relative margins are the load-bearing numbers. A
survivorship-bias-free vendor is the single highest-value data upgrade a
successor project could make.

\paragraph{Overfitting and data snooping.} Ten alphas, one composite,
and repeated evaluation on one out-of-sample window invite snooping; the
four gates exist to police exactly this, and the hyperparameter search
was designed so the out-of-sample window could not influence selection
(\cref{sec:eq-hp}). The residual risks are stated rather than waved
away: the walk-forward advantage is significant only within the tuned
configuration ($n=5$ per arm) and is otherwise a directional claim; the
flat hyperparameter surface cuts both ways, as robustness and as a
warning that tuning gains here are mostly noise; and the energy
forecasting results rest on a single six-month window with untuned
hyperparameters. The two synthetic negative controls and the
shuffle-label control (\cref{sec:eq-controls,sec:en-negcontrol}) are the
strongest argument that the pipeline as a whole does not snoop.

\paragraph{Transaction costs and tradability.} No result in this report
is net of costs: there is no fill model, no spread, no market impact.
The equity Sharpe of 1.37 is a paper Sharpe on daily bars, and a
five-day-horizon long-short over fifty names would face real turnover
costs. The energy track makes the point more brutally --- its strategy
loses money \emph{before} any costs are applied
(\cref{sec:en-clip}), and the E13b lesson (rank quality does not imply
profitability when tails carry the money) would only sharpen under a
cost model. The platform's stated production boundary
(\cref{sec:pl-layers}) is consistent with this: nothing here is a
trading system.

\paragraph{Regime dependence and external validity.} All equity claims
rest on a single 5.5-year window of one market regime, at one horizon
($k=5$ days), on one universe; all energy forecasting claims on six
months of one calendar year, at one horizon ($k=24$ hours), on twenty
zones. The attention analysis itself argues against over-generalising:
the learned structure is temporally stationary
(\cref{sec:eq-attention}), so the project's own evidence says the model
found a stable statistical regularity, not a regime-adaptive mechanism.
Whether either result survives a different half-decade is untested.

\paragraph{Reproducibility.} What is reproducible is the distribution,
not the point: the same seed produces different results on different
devices, for identified mechanical reasons (\cref{sec:eq-variance}).
All headline numbers are therefore means over at least five seeds on a
fixed device, the test suite pins the leakage-critical logic, and the
archived per-run artifacts allow every table in this report to be
regenerated. Bit-exact cross-device replication is not claimed and, on
CUDA, not attainable.

\paragraph{The open gate.} The relational composite never beat the best
single island alpha --- an OOS Sharpe of roughly 1.4 against roughly 3
--- and this report resists the temptation to soften that bar
post hoc. Two softer readings are available and honestly labelled as
softer: the composite beats nine of its ten inputs, and it is nearly
uncorrelated with all of them (maximum 0.26), which for a practitioner
suggests portfolio value that the max-of-singles gate does not measure.
A marginal-contribution-to-a-multifactor-portfolio reading is the
natural next evaluation, planned but not executed within this project's
scope.
